\documentclass[11pt]{ctexart}
\usepackage{amsmath,amssymb,enumitem,geometry}
\geometry{margin=1in}
\setlist[enumerate]{leftmargin=*}
\title{Week 2 Problem Set: Algorithmic Thinking, Data Structures, and Proof Patterns}
\author{TCS Self-Study OS}
\date{Generated 2026-05-06}

\begin{document}
\maketitle

说明：每题标注 difficulty 和 required concepts。先写 specification，再写 proof strategy。不要把实验输出当作证明。

\section{Day 1: Algorithm Specifications, Sorting, and Comparison Lower Bounds}
\begin{enumerate}
  \item \textbf{Basic. Concepts: specification.} 写出 sorting problem 的 formal specification：input、output、precondition、postcondition，并说明 n 的含义。
  \item \textbf{Basic. Concepts: stability, in-place.} 给出一个 stable 但不 in-place 的排序例子，以及一个 in-place 但不 stable 的排序例子；说明性质不同。
  \item \textbf{Medium. Concepts: loop invariant, insertion sort.} 证明 insertion sort 的外层循环 invariant：第 i 轮开始时 A[0..i) 是原前 i 个元素的 sorted permutation。
  \item \textbf{Medium. Concepts: merge sort recurrence.} 设 merge sort 对 n 个元素执行 T(n)=2T(n/2)+cn 次操作，T(1)=d。用 recursion tree 给出 T(n)=O(n log n) 的证明。
  \item \textbf{Medium. Concepts: comparison lower bound.} 在 decision tree model 中说明为什么 comparison sorting 的 leaves 至少有 n! 个。
  \item \textbf{Challenge. Concepts: lower bound model.} 解释为什么 counting sort 不违反 comparison sorting 的 Omega(n log n) lower bound。
  \item \textbf{Challenge. Concepts: upper and lower bounds.} 已知 merge sort worst-case O(n log n)，comparison sorting lower bound Omega(n log n)。在 comparison model 中能推出什么 Theta 结论？
\end{enumerate}

\section{Day 2: Divide and Conquer, Recurrence Proofs, and Selection}
\begin{enumerate}
  \item \textbf{Basic. Concepts: recursive correctness.} 写出 binary search 的 precondition、postcondition 和 loop invariant。
  \item \textbf{Basic. Concepts: recurrence.} 解 T(n)=T(n/2)+1, T(1)=1 的 Theta bound，并说明变量。
  \item \textbf{Medium. Concepts: merge sort correctness.} 用 induction 证明 merge sort 输出 sorted permutation，允许引用 merge procedure 正确性。
  \item \textbf{Medium. Concepts: substitution method.} 证明 T(n)<=2T(n/2)+n, T(1)<=1 满足 T(n)=O(n log n)。
  \item \textbf{Medium. Concepts: selection.} 说明为什么 selection problem 不需要完整排序，并给出 O(n log n) baseline。
  \item \textbf{Challenge. Concepts: quickselect expected time.} 给出 quickselect expected linear time 的直觉：什么 pivot 叫 good，good pivot 的概率为什么是常数级？
  \item \textbf{Challenge. Concepts: design.} 设计一个 divide-and-conquer 算法求数组最大值和最小值，并写 recurrence。
\end{enumerate}

\section{Day 3: Data Structures, Heaps, Hashing, and Amortized Analysis}
\begin{enumerate}
  \item \textbf{Basic. Concepts: ADT interface.} 为 priority queue 写出 insert 和 extract-min 的 interface specification。
  \item \textbf{Basic. Concepts: heap invariant.} 给数组 [1,3,2,7,5,4] 判断是否满足 min-heap invariant。
  \item \textbf{Medium. Concepts: heap insert correctness.} 证明 binary heap insert 后仍满足 heap invariant。
  \item \textbf{Medium. Concepts: hashing expected time.} 若 chaining hash table 有 n 个 keys 和 m 个 buckets，load factor alpha=n/m。在 simple uniform hashing 下 successful search expected time 为什么是 O(1+alpha)？
  \item \textbf{Medium. Concepts: dynamic array aggregate.} 证明从空 dynamic array 开始做 m 次 append，总复制成本 O(m)。
  \item \textbf{Challenge. Concepts: potential method.} 给 dynamic array doubling 提出一个 potential function 并说明它为什么能支付 resize。
  \item \textbf{Challenge. Concepts: design.} 设计一个 dictionary interface，用 hashing 实现，并说明 worst-case 与 expected-case。
\end{enumerate}

\section{Day 4: Graph Algorithms I: BFS, DFS, Shortest Paths}
\begin{enumerate}
  \item \textbf{Basic. Concepts: graph representation.} 比较 adjacency list 与 adjacency matrix 的空间复杂度。
  \item \textbf{Basic. Concepts: BFS trace.} 对 path graph 1-2-3-4 从 1 做 BFS，写出 distances 和 parent tree。
  \item \textbf{Medium. Concepts: BFS correctness.} 证明 BFS 在 unweighted graph 中第一次发现 vertex v 时 dist[v] 是 shortest distance。
  \item \textbf{Medium. Concepts: DFS not shortest.} 给出 DFS 找到非 shortest path 的反例。
  \item \textbf{Medium. Concepts: connected components.} 说明如何用 DFS/BFS 找 undirected graph 的 connected components，并分析复杂度。
  \item \textbf{Challenge. Concepts: topological ordering.} 证明 DAG 按 DFS decreasing finish time 给出 topological order。
  \item \textbf{Challenge. Concepts: Dijkstra limitation.} 给出 Dijkstra 在负权边上失败的原因或反例。
\end{enumerate}

\section{Day 5: Greedy Algorithms and Exchange Arguments}
\begin{enumerate}
  \item \textbf{Basic. Concepts: greedy counterexample.} 给 coin system {1,3,4} 和 amount 6，说明 largest-coin greedy 失败。
  \item \textbf{Basic. Concepts: interval scheduling trace.} 给 intervals [0,3),[1,2),[2,4),[3,5)，按 earliest finish greedy 选择哪些？
  \item \textbf{Medium. Concepts: exchange argument.} 写出 interval scheduling 第一个 greedy choice 的 exchange step。
  \item \textbf{Medium. Concepts: stays-ahead.} 对 interval scheduling 说明 stays-ahead 证明想证明什么不变量。
  \item \textbf{Medium. Concepts: fractional knapsack.} 证明 fractional knapsack density greedy 的安全性。
  \item \textbf{Challenge. Concepts: MST cut property.} 用 cut property 解释 Kruskal 为什么可以加入跨越两个 components 的最轻边。
  \item \textbf{Challenge. Concepts: greedy design.} 为“选择最多互不重叠 intervals”比较 earliest start 与 earliest finish，给 earliest start 反例。
\end{enumerate}

\section{Day 6: Dynamic Programming and Optimal Substructure}
\begin{enumerate}
  \item \textbf{Basic. Concepts: state definition.} 为 LIS 定义 dp[i]，说明它表示 ending at i 还是 prefix i。
  \item \textbf{Basic. Concepts: edit distance base cases.} 写出 edit distance dp[i][0] 与 dp[0][j]。
  \item \textbf{Medium. Concepts: LIS recurrence.} 写出 LIS O(n\textasciicircum{}2) recurrence 并证明覆盖所有情况。
  \item \textbf{Medium. Concepts: edit distance recurrence.} 写出 edit distance recurrence，并说明 match/substitute/delete/insert。
  \item \textbf{Medium. Concepts: computation order.} 为什么 edit distance tabulation 可按 i 从 0..n、j 从 0..m 填表？
  \item \textbf{Challenge. Concepts: knapsack recurrence.} 给 0/1 knapsack dp[i][w] recurrence 并说明 O(nW) 为什么是 pseudo-polynomial。
  \item \textbf{Challenge. Concepts: reconstruction.} 说明如何从 LIS dp 恢复一个实际 subsequence。
\end{enumerate}

\section{Day 7: Randomized Algorithms and Synthesis}
\begin{enumerate}
  \item \textbf{Basic. Concepts: Las Vegas vs Monte Carlo.} 区分 Las Vegas 和 Monte Carlo，并各给一个 Week 2 相关例子。
  \item \textbf{Basic. Concepts: sample space.} 为 randomized quickselect 定义 sample space。
  \item \textbf{Medium. Concepts: expected time.} 解释 expected running time 与 worst-case running time 的区别。
  \item \textbf{Medium. Concepts: failure probability.} 若 Monte Carlo test 每次独立失败概率 <=1/3，重复 r 次且全部失败才整体失败，给出 failure bound。
  \item \textbf{Medium. Concepts: amplification.} 若 majority vote amplification 使用 r 次 independent runs，整体错误事件是什么？
  \item \textbf{Challenge. Concepts: randomized quicksort.} 说明 randomized quicksort correctness 为什么不依赖 pivot distribution。
  \item \textbf{Challenge. Concepts: hashing probability.} 在 simple uniform hashing 中，两个不同 keys collision 的概率是多少？如何用于 expected bucket length？
\end{enumerate}

\end{document}
